\scanpage{7}
\diagram{\begin{tikzcd}
H^i({}_pT)\arrow[r]&{}_pH^i(T)\arrow[r,phantom,"\unsure{-|\to}" description]&0
\end{tikzcd}}
\diagram{\begin{tikzcd}
0\arrow[r,phantom,"\unsure{-|\to}" description]&H^i(T)_p\arrow[r,phantom,"\unsure{-|\to}" description]&H^{i+1}({}_pT)
\end{tikzcd}}
on voit que les conditions sont nécessaires.
Pour la suffisance, on tire d'abord le fait
que $H^{\unsure{i}}(G)\deleted{\unsure{(p)}}=\quad$ pour $i\geq n+2$ ; reste
à vérifier $H^{n+1}(G)=0$ donc que
$H^n(T)\deleted{\unsure{(p)}}\longrightarrow H^n(T')\deleted{\unsure{(p)}}$ est surjectif.
\deleted{\unsure{Si} $n\unsure{\ne}0$
\unsure{de tte} \unsure{façon} le \unsure{conoyau} est un \unsure{gr.}
de $p$-torsion, donc on peut se borner
à regarder
\diagram{\begin{tikzcd}
\deleted{H^n(T)(p)}\arrow[r]&\deleted{H^n(T')(p)}
\end{tikzcd}}
C'est un hom. de gr. $p$-divisibles,}

Noter que si $k'/k$ est une extension finie
séparable de $k$, alors pour tt $G'$ sur $k'$
on a
\[
H^i(k,\prod_{k'/k}G')\simeq H^i(k,G),
\]
ce qui prouve que l'hypothèse faite sur
les tores sur $k$ est \uline{\unsure{st.ble}} par
extension de base séparable. Soit $k'$ une
telle extension de degré premier à $p$, alors
l'hom. \unsure{can.}
\diagram{\begin{tikzcd}
\prod_{k'/k}G_{k'}\arrow[r]&G
\end{tikzcd}}
est \deleted{surjectif} \inserted{un \uline{\unsure{épim.}}}
\mspage{6}

